\begin{abstract}

This thesis presents a micro aerial vehicle that travels across water in a
sequence of hops, driven by hydrofoils that rotate with the airframe. Inspired
by stone-skipping, the vehicle descends onto a water surface with its lift
rotors switched off; the spinning hydrofoils submerge briefly, generate a large
upward impulse, and eject the vehicle back into the air, where it recovers
altitude under closed-loop control and repeats the cycle. Because water is
nearly three orders of magnitude denser than air, the same foils that produce a
negligible force in flight produce a large one on contact, so they cost the
vehicle little while airborne yet supply the whole of the impulse that returns
it to the air.

A quasi-steady blade-element model of the submersion phase is developed, giving
the vertical thrust and the resisting torque generated by a foil as it sweeps
through the water, and coupling them to the vertical and rotational motion of
the vehicle. Simulation across the design space identifies the entry velocity
and the foil inclination as the quantities that govern the hop, shows that the
radial extent of the foil is nearly inert once large enough, and exposes the
central tension of the manoeuvre: the rotation that generates the ejecting
force is itself consumed by the water.

The vehicle built to test this is a $90\;\mathrm{g}$ four-rotor airframe that
spins continuously about its vertical axis. Two rotors are mounted
tangentially and drive the rotation directly at a commanded thrust setpoint,
while two remain vertical and provide lift and attitude authority; because the
yaw drive is decoupled from the lift channel, the spin rate becomes a design
parameter rather than an equilibrium the airframe settles into, and the
rotation persists when the lift command is removed entirely. Flight control
for the resulting fast-rotating vehicle is formulated on its gyroscopic reduced
attitude dynamics, together with an online estimator that separates the
spin-synchronous measurement artefact from the attitude estimate.

Flight experiments in a motion-capture environment demonstrate the full cycle.
The descent is flown with the lift rotors commanded to zero and is arrested
with them still off, so the reversal is the work of the hydrofoils alone;
across eight contacts the vehicle enters the water at $-1.53$ to
$-1.85\;\mathrm{m\,s^{-1}}$ and leaves it climbing, recovering on average
$60\%$ of its entry speed. Consecutive hops are demonstrated with the vehicle
returning to within $8\;\mathrm{mm}$ of its previous release height, showing
that the cycle closes rather than that a single impact can be survived. The
vehicle remained controllable at spin rates approaching
$10{,}000^{\circ}\,\mathrm{s^{-1}}$, with no upper limit encountered.

This work establishes water-skipping as a usable mode of locomotion for small
aerial robots, and provides the model, the airframe, and the control needed to
pursue it further.

\end{abstract}
